\section{An exact rank criterion}\label{sec:rank}
The determinant comparison requires the matrix $V$ to have full row rank.
The necessary dimension inequality $nJ\geq M$ does not by itself establish
this. We use the same reduction at an inert prime as in the earlier bound,
but retain the individual degree and multiplicity estimates throughout the
argument. This gives a stronger sufficient criterion at the parameters needed
here.

The obstruction to direct reduction is the identity
$\bar z=z^p$ modulo an inert prime $p$. Thus $z^p-w$ vanishes at every
reduced node, and some of the rows become dependent. Dividing out known
powers of $p$ before reducing preserves more information. A relation among
the resulting rows would produce a polynomial with prescribed vanishing
along several curves. We will construct from it a nonzero polynomial in
one variable, and show that its zeros have total multiplicity greater than
its degree.

\begin{theorem}[Exact rank criterion]\label{thm:exact-rank}
Let $K\geq2$ and $S\geq1$ be integers, put $R=KS$, and let $1\leq J\leq R$.
Suppose that $p$ is inert in $\OO_D$, that
\[
 K<p<R\leq p^2,
\]
and that the nodes $z_1,\ldots,z_n\in\OO_D$ have distinct residues modulo
$p\OO_D$.
Set $d=\lfloor(R-1)/p\rfloor$, and define
\begin{equation}\label{eq:rank-caps}
 \begin{split}
 D_m&=p\left\lfloor\frac{R-mp-1}{K}\right\rfloor
             +\bigl((R-mp-1)\bmod K\bigr)\qquad(0\leq m\leq d),\\
 B_e&=\sum_{m=0}^e(D_m-m)\qquad(0\leq e\leq d).
 \end{split}
\end{equation}
For $l\geq0$ put $L_l=\min(p,\max(0,J-lp))$, and let
\begin{equation}\label{eq:rank-multiplicity}
 \mu_e=\sum_{l=0}^e\max(0,L_l-e+l)\qquad(0\leq e\leq d).
\end{equation}
If
\begin{equation}\label{eq:exact-rank-ineq}
 n\mu_e>B_e\qquad(0\leq e\leq d),
\end{equation}
then the matrix \eqref{eq:jet-matrix} formed with these nodes has full
row rank.
\end{theorem}

Here $D_m$ is the largest degree of a coefficient polynomial arising from
layer $m$ of the reduced relation. The quantity $B_e$ bounds the degree of
the one-variable polynomial constructed from a relation of degree $e$ in
the second variable, while $\mu_e$ bounds its order of vanishing at each
node. Keeping these quantities separately for every $e$ avoids the losses
in the earlier sufficient criterion.

\subsection{A relation after reduction modulo an inert prime}
Put $g=z^p-w$. The polynomials
\begin{equation}\label{eq:g-basis}
 g^m z^a w^s,\qquad
 0\leq m\leq d,\quad 0\leq a<p,\quad a+Ks<R-mp,
\end{equation}
form an integral basis of the space spanned by the triangular support
$z^rw^s$, $r+Ks<R$. Indeed, writing $r=mp+a$ identifies their leading
monomials with the original basis. Every other term replaces a factor
$z^p$ by $-w$, decreasing the weighted degree by $p-K>0$.
The change of basis is therefore unitriangular.

At node $i$ write
\[
 g(z_i+t,\bar z_i)=pU_i(t)+t^p,\qquad
 U_i(t)=\frac{(z_i+t)^p-t^p-\bar z_i}{p}\in\OO_D[t].
\]
The constant coefficient is integral because $p$ is inert, and the remaining
coefficients are integral because $p$ divides $\binom pv$ for $0<v<p$.
Divide each row in layer $m$ by $p^m$, and multiply each column of derivative
order $j=lp+v$, $0\leq v<p$, by $p^l$. Denote the resulting matrix by
$\widetilde V$. Its entries are still integral: a term using $t^{pq}$ in
the expansion of $(pU_i+t^p)^m$ can contribute to order $j$ only if $q\leq l$,
and its coefficient after rescaling contains $p^{l-q}$.
Only $q=l$ can survive reduction modulo $p$.

Write $k=\OO_D/p\OO_D$, which is a field, and put
$\xi_i=z_i\bmod p$ and $u_i(z)=(U_i\bmod p)(z-\xi_i)$.
For $l\leq m$ the reduced entry is
\begin{equation}\label{eq:scaled-entry}
 \binom ml[t^v]\bigl(u_i(\xi_i+t)^{m-l}(\xi_i+t)^a\xi_i^{ps}\bigr);
\end{equation}
for $l>m$ it is zero. The basis change and the scalings are invertible in
characteristic zero. Consequently, full row rank of $\widetilde V\bmod p$
will imply full row rank of $V$. These scalings are used only for this
purpose; the size and divisibility estimates concern the original matrix.

Suppose that the reduced rows satisfy a nontrivial relation, with
coefficients $c_{m,a,s}\in k$. Define
\begin{equation}\label{eq:rank-F}
 F(z,u)=\sum_{m=0}^d H_m(z)u^m,\qquad
 H_m(z)=\sum_{a,s}c_{m,a,s}z^{a+ps}.
\end{equation}
The exponents $a+ps$ within a layer are distinct because $0\leq a<p$.
Thus $F\ne0$. Let $e=\deg_u F\leq d<p$, where the last inequality follows from
$R\leq p^2$.
For $t=R-mp-1$, the maximum of $a+ps$ subject to
$a+Ks\leq t$ and $0\leq a<p$ is attained at
$s=\lfloor t/K\rfloor$, $a=t\bmod K$, since $p>K$.
It follows that
\begin{equation}\label{eq:rank-degree}
 \deg H_m\leq D_m.
\end{equation}
We will also use $D_m-D_{m+1}\geq p$ for $m<d$.
To see this without manipulating the floors, take a monomial attaining
$D_{m+1}$ and increase $s$ by one. Its weighted degree increases by $K<p$,
so it is admissible in layer $m$, and its degree $a+ps$ increases by $p$.

There are $L_l$ derivative orders $j=lp+v$ with $0\leq j<J$.
Since $(\xi_i+t)^{ps}\equiv\xi_i^{ps}\pmod{t^p}$ in characteristic $p$,
we may replace $\xi_i^{ps}$ in \eqref{eq:scaled-entry} by
$(\xi_i+t)^{ps}$ when extracting coefficients of order less than $L_l$.
The row relation therefore gives
\begin{equation}\label{eq:local-hasse-vanish}
 \ord_{z=\xi_i}\bigl((\partial_u^{[l]}F)(z,u_i(z))\bigr)\geq L_l
 \qquad(1\leq i\leq n,\ 0\leq l\leq e).
\end{equation}
Here $\partial_u^{[l]}$ is the Hasse derivative, defined by
$\partial_u^{[l]}u^m=\binom ml u^{m-l}$.

\subsection{A nonzero polynomial of controlled degree}
Differentiating the identity defining $U_i$ before reducing modulo $p$ gives
\begin{equation}\label{eq:ui-prime}
 u_i'(z)=z^{p-1}-(z-\xi_i)^{p-1}.
\end{equation}
Thus differentiation along the different curves $u=u_i(z)$ agrees, up to
a term vanishing to order $p-1$ at the relevant node, with the common operator
\[
 \mathcal D=\partial_z+z^{p-1}\partial_u.
\]
On $k(z)[u]$ this operator satisfies
\begin{equation}\label{eq:Dp-identity}
 \mathcal D^p=-\partial_u.
\end{equation}
Indeed, the $p$th iterate of a derivation in characteristic $p$ is a
derivation. It sends $z$ to zero and $u$ to $(p-1)!=-1$, which determines
its action on $k(z)[u]$.

We claim that $F,\mathcal DF,\ldots,\mathcal D^eF$ are linearly independent
over $k(z)$. Otherwise, since $F\ne0$, let $1\leq r\leq e$ be the first index at
which a dependence occurs. The span of $F,\ldots,\mathcal D^{r-1}F$ is then stable
under $\mathcal D$, including when $\mathcal D$ differentiates rational
coefficients. By \eqref{eq:Dp-identity} it is also stable under $\partial_u$.
It must therefore contain $F,\partial_uF,\ldots,\partial_u^eF$.
These are independent: since $e<p$, they have distinct degrees
$e,e-1,\ldots,0$ in $u$, with nonzero leading coefficients.
This contradicts the dimension $r\leq e$.

Consequently the coefficient determinant
\begin{equation}\label{eq:WF-def}
 W_F(z)=\det\bigl([u^m]\mathcal D^jF\bigr)_{0\leq j,m\leq e}
\end{equation}
is a nonzero polynomial. The identity
\[
 [u^m]\mathcal DG=([u^m]G)' +(m+1)z^{p-1}[u^{m+1}]G
\]
and $D_m-D_{m+1}\geq p$ give, by induction on $j$,
\[
 \deg\bigl([u^m]\mathcal D^jF\bigr)\leq D_m-j.
\]
Terms with $m>e$ are zero, and a negative degree bound means the
corresponding entry is zero. Each determinant term uses every row and every
column exactly once, so
\begin{equation}\label{eq:WF-degree}
 \deg W_F\leq\sum_{m=0}^eD_m-\sum_{j=0}^ej=B_e.
\end{equation}

\subsection{The multiplicity at each node}
Fix $i$. Expressing each row polynomial of \eqref{eq:WF-def} in powers of
$u-u_i(z)$ makes a unitriangular change of columns. It preserves the
determinant and replaces its entries by
\[
 T^{(i)}_{j,l}(z)=(\partial_u^{[l]}\mathcal D^jF)(z,u_i(z)).
\]
We first show that
\begin{equation}\label{eq:local-entry-order}
 \ord_{\xi_i}T^{(i)}_{j,l}\geq\max(0,L_l-j).
\end{equation}
Let $\mathcal D_i=\partial_z+u_i'(z)\partial_u$.
By \eqref{eq:ui-prime},
$\mathcal D-\mathcal D_i=(z-\xi_i)^{p-1}\partial_u$.
Repeated application lowers the order of this error by at most one at each
step; hence, for every polynomial $G$ and $0\leq j<p$,
\[
 \mathcal D^jG-\mathcal D_i^jG
       \in(z-\xi_i)^{p-j}k[z,u].
\]
Both operators commute with the Hasse derivatives in $u$.
After substituting $u=u_i(z)$, the entry $T^{(i)}_{j,l}$ therefore agrees,
modulo $(z-\xi_i)^{p-j}$, with the $j$th ordinary derivative of
$(\partial_u^{[l]}F)(z,u_i(z))$.
Equation~\eqref{eq:local-hasse-vanish} bounds its order from below by $L_l-j$ when
this is positive. Since $L_l\leq p$, the error has at least the same
order, proving \eqref{eq:local-entry-order}.

To use all these bounds together, consider a term in the determinant
expansion. Its order is at least
$\sum_{l=0}^e(L_l-j_l)_+$, where the $j_l$ are a permutation of
$0,\ldots,e$. For $a\geq b$ and $x\geq y$ one has
\[
 (a-x)_++(b-y)_+\leq(a-y)_++(b-x)_+.
\]
Thus pairing the largest $L_l$ with the largest row index minimizes this
sum. Since $L_0\geq\cdots\geq L_e$, its minimum is
$\sum_{l=0}^e(L_l-e+l)_+=\mu_e$.
Every determinant term has at least this order, and so does their sum:
\[
 \ord_{\xi_i}W_F\geq\mu_e.
\]
Since $W_F$ is nonzero and the residues are distinct, \eqref{eq:WF-degree} gives
\begingroup\postdisplaypenalty=10000
\[
 n\mu_e\leq\deg W_F\leq B_e,
\]
\endgroup
contrary to \eqref{eq:exact-rank-ineq}. This proves Theorem~\ref{thm:exact-rank}.
